% Source: https://artofproblemsolving.com/community/c3h2369551_counting_help
% Stable identity: aops-topic-2369551-q01
Twenty-five students sit around a round table. Three of them are chosen, all choices of three students being equally likely.
\begin{enumerate}
\item[(a)] Find the probability that at least two of the three chosen students were sitting next to each other.
\item[(b)] A posted solution lets $a,b,c$ be the numbers of unchosen students in the three gaps between consecutive chosen students (going around the table), so that $a+b+c=22$. It says that $a+b+c=22$ has $\dbinom{21}{2}$ solutions with $a,b,c\ge1$ and $\dbinom{24}{2}$ solutions with $a,b,c\ge0$, and that the probability that no two chosen students are neighbours is $\dfrac{\dbinom{21}{2}}{\dbinom{24}{2}}$. Explain where $\dbinom{21}{2}$ and $\dbinom{24}{2}$ come from, and why this ratio is correct even though a choice of three students does not correspond to just one triple $(a,b,c)$.
\end{enumerate}
